\DocumentMetadata{
  pdfversion=2.0,
  pdfstandard=ua-2,
  lang=en-US,
  tagging=on,
  tagging-setup={math/setup=mathml-SE,tabsorder=structure}
}
\documentclass[11pt,letterpaper]{article}
\usepackage[margin=0.68in,includefoot]{geometry}
\usepackage{newcomputermodern}
\usepackage{amsmath,amscd,mathtools}
\usepackage{tikz,adjustbox}
\providecommand{\square}{\mdlgwhtsquare}
\usepackage[hidelinks]{hyperref}
\hypersetup{
  pdftitle={Algebra First-Year Exam, January 2018, Part 2},
  pdfauthor={University of Florida Department of Mathematics},
  pdfsubject={Graduate mathematics examination},
  pdfdisplaydoctitle=true,
  bookmarksopen=true
}
\setcounter{secnumdepth}{0}
\setlength{\parindent}{0pt}
\setlength{\parskip}{0.28em}
\setlength{\emergencystretch}{3em}
\setlength{\leftmargini}{2.15em}
\setlength{\leftmarginii}{2.65em}
\setlength{\leftmarginiii}{3em}
\renewcommand{\labelenumi}{\textbf{\theenumi.}}
\pagestyle{empty}
\raggedbottom
\allowdisplaybreaks
\newenvironment{examproblems}[1]{%
  \begin{enumerate}%
  \setcounter{enumi}{#1}%
  \setlength{\topsep}{0.35em}%
  \setlength{\partopsep}{0pt}%
  \setlength{\itemsep}{0.48em}%
  \setlength{\parsep}{0.18em}%
}{\end{enumerate}}
\newenvironment{examsubparts}{%
  \begin{enumerate}%
  \setlength{\topsep}{0.18em}%
  \setlength{\partopsep}{0pt}%
  \setlength{\itemsep}{0.16em}%
  \setlength{\parsep}{0.1em}%
}{\end{enumerate}}
\newenvironment{examinstructions}{%
  \par\begingroup\itshape
}{%
  \par\endgroup\vspace{0.18em}
}
\makeatletter
\renewcommand\section{\@startsection{section}{1}{0pt}%
  {-1.25ex plus -.25ex minus -.1ex}%
  {0.55ex plus .1ex}%
  {\normalfont\large\bfseries}}
\renewcommand{\@maketitle}{%
  \newpage\null\vskip -1em%
  {\footnotesize\bfseries
    \href{https://www.ufl.edu/}{University of Florida}, \href{https://math.ufl.edu/}{Department of Mathematics}\par}%
  \vskip -0.5em%
  \tagstructbegin{tag=Title}%
    {\LARGE\bfseries\href{https://gma.math.ufl.edu/exams/algebra-first-year/}{Algebra First-Year Exam}, January 2018, Part 2\par}%
  \tagstructend%
  \vskip 0.25em%
  \tagmcbegin{artifact}%
    \hrule height 1pt%
  \tagmcend%
  \vskip 1em%
}
\makeatother
\title{Algebra First-Year Exam, January 2018, Part 2}
\author{}
\date{}
\begin{document}
\maketitle
\thispagestyle{empty}
\begin{examinstructions}
Answer four problems. Write your answers clearly in complete English sentences. You may quote results (within reason) as long as you state them clearly.
\end{examinstructions}
\begin{examproblems}{0}
\item
Show that the ring 
\[
\mathbb{Z}[\omega]=\{a+b\omega\in\mathbb{C}\mid a,b\in\mathbb{Z}\},\qquad \omega=\frac{1+i\sqrt{3}}{2}
\]
 is a Euclidean domain.
\item
Let~\(R\) be a commutative ring with identity. A subset \(S\subseteq R\) is multiplicative if it is nonempty and is closed under multiplication. Show that if~\(S\) is a multiplicative subset of~\(R\) not containing~\(0\), there is a prime ideal \(\mathfrak{p}\subset R\) such that \(\mathfrak{p}\cap S=\varnothing\). (First show using Zorn’s lemma that the set of ideals \(I\subset R\) such that \(I\cap S=\varnothing\) has a maximal element. Then show, using an argument by contradiction that any such maximal element is prime).
\item
Prove the following form of Eisenstein’s criterion: Let~\(R\) be an integral domain and \(P\subset R\) a prime ideal. Suppose 
\[
f(X)=X^d+a_{d-1}X^{d-1}+\cdots+a_1X+a_0\in R[X]
\]
 is a monic polynomial such that \(a_i\in P\) for \(i<d\) and \(a_0\notin P^2\). Then \(f(X)\) is irreducible in \(R[X]\).
\item
Let~\(R\) be an integral domain.
\begin{examsubparts}
\item[\textnormal{(i)}]
Define what it means for an element of~\(R\) to be irreducible.
\item[\textnormal{(ii)}]
Define what it means for an element of~\(R\) to be prime.
\item[\textnormal{(iii)}]
Show that a prime element is irreducible.
\end{examsubparts}
\item
Suppose~\(R\) is a commutative ring with identity, \(M\) and~\(N\) are~\(R\)-modules and \(f:M\to N\) is an~\(R\)-module homomorphism. Let \(I\subset R\) be an ideal.
\begin{examsubparts}
\item[\textnormal{(i)}]
Show that there is a homomorphism \(\overline{f}:M/IM\to N/IN\) such that 
\[
\overline{f}(m+IM)=f(m)+IN
\]
 for all \(m\in M\).
\item[\textnormal{(ii)}]
Show that if~\(I\) is nilpotent (i.e. \(I^k=0\) for some \(k>0\)) and \(\overline{f}\) is surjective then~\(f\) is surjective.
\end{examsubparts}
\item
Find representatives of every similarity class of \(6\times 6\) matrices with rational coefficients whose minimal polynomial is \((x+1)^2(x-1)\).
\end{examproblems}
\end{document}
